\chapter*{General Conclusion}
\addcontentsline{toc}{chapter}{General Conclusion}
\markboth{General Conclusion}{General Conclusion}

This internship addressed a real problem: release preparation at Forvia
Informatique Tunisie was slow and spread across multiple Azure DevOps screens
with no unified view of how artifacts related to one another. The work covered
the full development cycle, from architecture and implementation to integration,
validation, and deployment.

The result is a modular traceability platform. The backend handles recursive
graph traversal across work items, pull requests, commits, and tags, with cycle
prevention, batch processing, and SSE streaming for progressive feedback. The
frontend gives operators a single flow from work item selection to
promotion-ready summaries, replacing a process that previously required
navigating several disconnected pages. A lightweight LLM review feature was
also integrated at commit level, keeping human validation as the final step.

Beyond the delivered features, the project confirmed a few practical lessons.
Incremental module delivery kept integration risks low. Strict DTO contracts
between backend and frontend prevented most regressions. Logging and structured
error handling proved essential once flows became multi-step and harder to
trace manually.

The platform is functional and deployed, but also designed to evolve. The
traceability graph it builds is machine-readable, which makes it a natural
foundation for a future planning assistant able to recommend repository change
paths and target versions directly from a work item or user story. That next
step would take MES-X from dependency exploration to guided release preparation,
which is the logical continuation of what was built here.